% !TeX root = ../apunte_mecanica_clasica.tex
\IfFileExists{../apunte_mecanica_clasica.tex}{%
  \documentclass[../apunte_mecanica_clasica.tex]{subfiles}%
}{%
  \documentclass[apunte_mecanica_clasica.tex]{subfiles}%
}
\begin{document}

\chapter{Mecánica}
% Figuras extraídas desde Clase MC 2-1.pdf y organizadas en figures/08_mecanica/.

\section{Movimiento, fuerza e interacción}

\subsection{Primera ley de Newton}

La \textbf{primera ley de Newton} establece:
\begin{quote}
Todo cuerpo continúa en su estado de reposo o movimiento uniforme en línea recta, a menos que una fuerza externa neta actúe sobre él.
\end{quote}

La primera ley de Newton establece, en primer lugar, el movimiento como la piedra fundamental de la física, de acuerdo con las ideas de Aristóteles.

Segundo; introduce el movimiento con velocidad constante como un movimiento patrón, así como el metro, el kilogramo, etc., que nos permite detectar la presencia de fuerzas. Esto es, si se mueve en línea recta con velocidad constante entonces no hay fuerza neta, de lo contrario si hay fuerza neta. Así, dado un movimiento de un cuerpo lo comparamos con el movimiento patrón para detectar fuerzas.

Tercero; introduce el concepto \textbf{fuerza}:
\begin{quote}
Magnitud física que modifica el estado de movimiento de un cuerpo (lo acelera, frena, cambia el sentido, etc.) o bien lo deforma.
\end{quote}

\subsection{Interacción y modelo de fuerza}

Para entender lo que es el concepto fuerza debemos introducir el concepto \textbf{interacción}:
\begin{quote}
Es la acción de una parte del Universo sobre otra parte del Universo.
\end{quote}

La palabra acción indica que una persona, animal o cosa (material o inmaterial) está haciendo algo sobre otra persona, animal o cosa, o sea está actuando (de manera voluntaria o involuntaria, de pensamiento, palabra u obra) sobre algo, lo que normalmente implica movimiento o cambio de estado o situación y afecta o influye en una persona, animal o cosa.

En mecánica nos interesa la acción que genera movimiento; una parte 1 del Universo (agente externo) al actuar sobre otra parte 2 del Universo (sistema) genera sobre 2 un cambio en su estado de movimiento. Por ejemplo, una estrella interactuando sobre un planeta, una bola de billar en movimiento interactuando sobre otra bola de billar en reposo, etc.

Para desarrollar la mecánica debemos introducir una magnitud física que modele la interacción, llamada \textbf{fuerza}:
\begin{quote}
La \textbf{fuerza} es el modelo físico-matemático de una interacción.
\end{quote}

En símbolos,
\begin{equation*}
\text{Interacción} \longrightarrow \vec F .
\end{equation*}
Las fuerzas se representan mediante vectores, ya que poseen magnitud y dirección.

\subsection{Fuerzas inerciales y observadores inerciales}

Newton, como se mencionó, introduce el movimiento como la piedra fundamental de la mecánica y no la fuerza. La razón es que un observador que estudia un sistema puede observar fuerzas actuando sobre el sistema. Sin embargo, algunas de estas fuerzas no provienen de una interacción, no hay una porción del Universo actuando sobre el sistema. Por ejemplo, un observador que está girando y siente una fuerza que lo tira hacia afuera, esta fuerza es debido a la inercia y no la acción de otra porción del Universo. Esta fuerza se llama fuerza centrífuga. A estas fuerzas Newton las llamó fuerzas ficticias.

Por lo tanto, se concluye que en la teoría de Newton una interacción implica una fuerza, pero una fuerza no implica una interacción:
\begin{equation*}
\text{Interacción} \Longleftarrow \vec F.
\end{equation*}

Entonces aparece la pregunta: ¿cómo sabremos si las fuerzas que están presentes son ficticias o no? Las fuerzas ficticias están presentes cuando el observador está acelerado, porque actúa la inercia que trata de llevar al observador acelerado a un movimiento con velocidad constante, por lo tanto aparecerán las fuerzas ficticias.

En consecuencia, los observadores que no medirán fuerzas ficticias son los llamados \textbf{observadores inerciales}:
\begin{quote}
Son observadores que se mueven con velocidades relativas constantes entre ellos y el espacio absoluto.
\end{quote}
Por lo tanto, en la teoría de Newton hay observadores privilegiados.

Muchos físicos contemporáneos y posteriores a Newton no estuvieron de acuerdo con el nombre de fuerzas ficticias, porque después de todo son reales y tienen efecto sobre los cuerpos. Por ejemplo, viajar en un auto a gran velocidad y frenar repentinamente, si no se tiene puesto el cinturón de seguridad el conductor, y pasajeros si los hay, se proyectarán hacia adelante, lo que puede ser peligroso. Por lo tanto, los físicos le cambiaron el nombre de fuerzas ficticias a fuerzas inerciales. Físicos posteriores las llamaron fuerzas universales, puesto que estas fuerzas son proporcionales a la masa.

\section{Espacio absoluto y movimiento absoluto}

\subsection{Espacio absoluto}

La ley de la inercia en la teoría de Newton no tiene un agente externo que represente una interacción sobre el sistema, esto es no hay un agente externo que interactúe sobre el sistema que genere una fuerza inercial. La inercia sería un efecto no causal.

Para resolver este problema introduce el llamado espacio absoluto:
\begin{quote}
El espacio absoluto es la referencia absoluta del movimiento, es uniforme, lo que implica que todas sus partes son idénticas y existe independiente de los objetos que pueblan el Universo.
\end{quote}

Este espacio absoluto actúa sobre los cuerpos que están acelerados a través de las fuerzas inerciales, fuerzas que tratan de llevar los cuerpos acelerados a movimientos con velocidad constantes. El espacio absoluto actúa sobre todos los cuerpos del Universo.

\subsection{Espacio como entidad o relación}

En la cultura occidental existen dos maneras tradicionales de entender lo que es el espacio: como una \textbf{entidad} o una \textbf{relación}.

El espacio es una entidad: esto significa que el espacio existe aún cuando no haya nada más que el espacio. El espacio existe por sí mismo, y los objetos pueden moverse en él. Esta es la manera que Newton describe el espacio, y es conocido como espacio absoluto. Esta es también la manera en que el espacio-tiempo es entendido en relatividad especial. Aunque considerada desde tiempos antiguos (en la tradición de Demócrito), esta manera de entender el espacio no era la visión tradicional dominante en la cultura occidental. La visión dominante desde Aristóteles a Descartes fue entender el espacio como una relación.

El espacio es una relación: esto significa que el mundo está hecho por objetos físicos, o entidades físicas. Estos objetos tiene la propiedad que ellos pueden o no estar en contacto con otros. El espacio es este contacto, o contigüidad, o relación de adyacencia entre objetos. Aristóteles define la localización espacial de un objeto como el conjunto de objetos que lo rodean. Esto es, el espacio relacional.

\subsection{Movimiento absoluto y movimiento relativo}

Estrictamente conectada a estas dos maneras de entender el espacio, hay dos maneras de entender el movimiento.

\textbf{Movimiento absoluto}: si el espacio es una entidad, el movimiento puede ser definido como el viaje desde una parte del espacio a otra parte del espacio, es decir, como el cambio de posición desde un punto a otro del espacio. Esta es la manera como Newton define el movimiento.

\textbf{Movimiento relativo}: si el espacio es una relación, el movimiento puede sólo ser definido como el viaje desde la contigüidad de un objeto a la contigüidad de otro objeto.

La segunda visión es como Aristóteles y Descartes definen el movimiento, en palabras de Descartes (\emph{Principia Philosophiae}):
\begin{quote}
Podemos decir que el movimiento es la transferencia de una parte de la materia o de un cuerpo, desde la vecindad de aquellos cuerpos inmediatamente contiguos a él, y considerando el resto, en la vecindad de algún otro.
\end{quote}

Aristóteles postuló que el movimiento es relativo, y lo ilustra con el ejemplo de un hombre que camina sobre un bote. El hombre se mueve con respecto al bote, el cual se mueve con respecto del agua del río, la que se mueve con respecto de la arena, etc.

El relacionismo de Aristóteles es respaldado por el hecho que existen objetos preferidos que pueden ser usados como referencias preferidas: la tierra en el centro del universo y las esferas celestes, en particular las estrellas fijas muy lejos. De modo que podemos decir alguna cosa que está en movimiento si se mueve con respecto a la tierra. Sin embargo, hay dos sistemas preferidos en la cosmología antigua: la tierra y las llamadas estrellas fijas, y las dos rotan una respecto de la otra. Los pensadores de la edad media no ocultaron este punto, y discutieron largamente si las estrellas rotan alrededor de la tierra o si es la tierra que rota bajo las estrellas. Debe ser mencionado que en el siglo XIV Burilan concluyó que ningún punto de vista es más verdadero que la otra. Oresme estudió la rotación de la tierra más de un siglo antes que Copérnico.

Para un físico, la cuestión es cual de estas dos formas de pensamiento acerca del espacio y del movimiento permite una descripción más efectiva de la naturaleza. Para Newton, el espacio es absoluto y el movimiento es absoluto. Esta es una segunda violación del Cartesianismo.

\subsection{Críticas al espacio absoluto}

Las críticas al espacio absoluto de Newton no se hicieron esperar:
\begin{itemize}
    \item No especifica donde está este espacio absoluto.
    \item No explica de que está construido este espacio absoluto.
    \item Este espacio actúa sobre nosotros, pero nosotros no podemos actuar sobre él.
\end{itemize}
El argumento más fuerte usado por Newton a favor de su postulado fue la evidencia empírica discutida en su famoso experimento del balde.

\newpage
\section{Experimento del balde de Newton}

\subsection{Descripción del experimento}

En 1689, Newton quiso demostrar que el espacio absoluto existía. Llenó un balde de agua, lo ató a una cuerda y lo colgó de una barra fija.

Primera etapa: Newton retorció la cuerda al máximo, y lo soltó. El balde comienza a girar cada vez más rápido con respecto al agua, que está en reposo con respecto al balde.

Segunda etapa: durante unos momentos el agua sigue en reposo. El balde gira con respecto al agua, la superficie del agua permanece plana. Pero pasados unos instantes, el movimiento del balde se transfiere también al agua. El agua comienza a girar, debido al roce balde-agua, formando un remolino y su superficie se vuelve cóncava, así el agua gira con respecto al balde y viceversa.

Tercera etapa: Newton frena el balde con las manos, pero el agua sigue cóncava y girando por unos instantes más con respecto al balde.

Cuarta etapa: cuando el agua frena, su superficie vuelve a ser plana, el balde y el agua no rotan uno con respecto al otro.

\begin{apuntebloque}{Resumen}
\begin{enumerate}
    \item El balde rota con respecto al agua, la superficie del agua es plana.
    \item Por roce, balde y agua rotan juntos, la superficie del agua es cóncava.
    \item El agua rota con respecto al balde que está en reposo, la superficie del agua es cóncava.
    \item El balde y el agua están en reposo, la superficie del agua es plana.
\end{enumerate}
\end{apuntebloque}

\subsection{Argumento físico de Newton}

Si el movimiento es un cambio de lugar con respecto a los objetos que lo rodean, como exige Descartes, entonces debemos decir que en (1) el agua rota (con respecto al balde, el cual la rodea), mientras que en (2) el agua está en reposo (con respecto al balde). Pero, hace notar Newton, la rotación produce aquí un efecto físico: la concavidad de la superficie del agua. Este efecto aparece en (2), y no en (1). El efecto concavidad aparece en (2) cuando el agua debiera estar en reposo con respecto al balde, porque giran junto, no cuando el agua se mueve con respecto del balde. Por lo tanto, la rotación que produce el efecto físico no es la rotación con respecto al balde. ¿Cuál es la rotación que produce el efecto físico?

El razonamiento de Newton fue el siguiente: ¿cómo podemos distinguir que algo se mueve? Sólo hay una forma de saberlo, y esta consiste en tomar como referencia algo que no se mueve. El movimiento sólo existe con respecto a algo. Con cualquier referencia espacial sucede lo mismo: no hay izquierda sin derecha, ni arriba sin abajo, ni rápido sin lento.

El agua no se mueve con respecto al balde, sus tiempos no coinciden. Cuando el balde ya se mueve, el agua aún no, y cuando el balde ya se ha parado, el agua sigue moviéndose y sabemos que se mueve porque su superficie cambia de forma, siente la fuerza del movimiento. Entonces ¿con respecto a qué se mueve?

Newton responde:
\begin{quote}
Es la rotación con respecto al espacio absoluto. La concavidad de la superficie del agua es un efecto del movimiento absoluto del agua: es el movimiento con respecto al espacio absoluto. Esto, afirma Newton, prueba la existencia del espacio absoluto.
\end{quote}

\subsection{Relatividad galileana y relacionismo}

Para entender la razón del porqué los argumentos de Newton se mantuvieron sin cambios durante tres siglos, es necesario clarificar algunos mal entendidos comunes. El relacionismo, a saber, es la idea de que el movimiento puede ser definido sólo en relación con otros objetos, no debería ser confundido con la relatividad de Galileo.

La relatividad de Galileo es:
\begin{quote}
Realizando experimentos mecánico, el movimiento rectilíneo y uniforme es indistinguible del reposo.
\end{quote}

Concluimos que la velocidad y solo la velocidad es relativa. El relacionismo, por otro lado, afirma que cualquier movimiento, incluido el movimiento en zigzag, es a priori indistinguible del reposo. La afirmación de Galileo de la relatividad asume una definición no relacional del movimiento. Entonces: ¿movimiento rectilíneo y uniforme con respecto a qué?

Cuando Newton afirmó que el movimiento con respecto al espacio absoluto es real y físico, omitió insistir que aun el movimiento rectilíneo y uniforme es absoluto. Esto causó un acalorado debate, debido a que no hay ningún efecto físico del movimiento inercial. En este sentido, para Newton, el movimiento inercial y la velocidad son conceptos relativos. En efecto, movimiento y velocidad son conceptos relativos en mecánica de Newton.

Lo que Newton necesitó para formular la dinámica (interacciones), no es saber si el movimiento inercial es o no relativo, sino lo que necesitaba es si el movimiento acelerado, ejemplificado por la rotación del agua en el balde, es relativo o absoluto. La pregunta aquí no es si existe o no un espacio absoluto con respecto al cual pueda ser definida la velocidad. La pregunta es si existe o no un espacio absoluto con respecto al cual pueda ser definida la aceleración.

\subsection{Aceleración absoluta y segunda ley}

Newton respondió, basado en el experimento del balde, que existe un espacio absoluto que permite definir la \textbf{aceleración}. Sin esta respuesta, la segunda ley de Newton
\begin{equation*}
\vec F = m\vec a
\end{equation*}
no tendría sentido.

La oposición al espacio absoluto de Newton fue aún más fuerte que la oposición a su fuerza de acción a distancia. Leibnitz y su escuela argumentaron, con gran vehemencia, en contra del movimiento absoluto de Newton y del uso de la aceleración absoluta. Las dudas no desaparecieron con el correr de los siglos y persistió el sentimiento que algo estaba malo en el argumento de Newton.

Ernest Mach retornó al problema sugiriendo que el argumento del balde de Newton podría estar equivocado a causa de que el agua no rota con respecto del espacio absoluto, sino que ella rota con respecto al contenido total de materia del universo. El pensamiento de Mach tuvo gran influencia en Einstein cuando formuló Relatividad General.

\section{Segunda ley de Newton y fuerza inercial}

La \textbf{segunda ley de Newton} se expresa en la forma
\begin{equation}
    m\vec a = \vec F .
    \label{eq:mecanica-segunda-newton}
\end{equation}
Esta ley unifica la cinemática (movimiento) con la dinámica (interacción), esto es
\begin{equation*}
    m \times \text{Cinemática} = \text{Dinámica}.
\end{equation*}
La unificación se llama Mecánica. La constante de proporcionalidad $m$ se llama masa inercial, y mide la flojera de la naturaleza, esto es la masa inercial mide la oposición a cambiar el estado de movimiento de un sistema.

A diferencia a cómo interpreta la matemática cada lado de una ecuación, en física; cada lado de la una ecuación tiene una interpretación física, en general, diferente, debido a que las ecuaciones físicas son balances entre sistemas, o sea intercambio físicos entre sistemas.

En mecánica, la dinámica (interacciones) son las fuerzas externas que actúan sobre un sistema generando movimiento (cinemática), o sea la dinámica es la fuente de la cinemática. Por lo tanto, un término al escribirse en un lado de la ecuación cambia su significado físico cuando se pasa al otro lado de la ecuación.

\subsection{Transformación entre observadores}

Por ejemplo, sean dos observadores que analizan el movimiento de una partícula $P$; $O$ en reposo y $O'$ un observador que se mueve con respecto a $O$, no uniformemente. Estos están relacionados a través de la expresión
\begin{equation}
    \vec x_{P/O}=\vec x_{P/O'}+\vec x_{O'/O}.
    \label{eq:mecanica-posicion-relativa-observadores}
\end{equation}

\begin{figure}[H]
    \centering
    \IfFileExists{figures/08_mecanica/observadores_no_inerciales.png}{%
        \includegraphics[width=0.55\linewidth]{figures/08_mecanica/observadores_no_inerciales.png}%
    }{%
        \fbox{\parbox{0.55\linewidth}{\centering Figura pendiente: observadores $O$, $O'$ y partícula $P$.}}%
    }
    \caption{Relación geométrica entre los observadores $O$, $O'$ y la partícula $P$.}
    \label{fig:mecanica-observadores-no-inerciales}
\end{figure}

La segunda ley de Newton~\eqref{eq:mecanica-segunda-newton} se escribe con respecto al observador $O$ como
\begin{equation}
    m\left.\frac{d^2\vec x}{dt^2}\right|_{O}=\vec F,
    \label{eq:mecanica-segunda-newton-observador-o}
\end{equation}
donde la fuerza $\vec F$ es absoluta. Sin embargo, si la escribimos con respecto a $O'$, reemplazando~\eqref{eq:mecanica-posicion-relativa-observadores} en~\eqref{eq:mecanica-segunda-newton-observador-o} se encuentra
\begin{equation}
    m\left(\frac{d^2\vec x_{P/O'}}{dt^2}+\frac{d^2\vec x_{O'/O}}{dt^2}\right)=\vec F.
    \label{eq:mecanica-segunda-newton-reemplazo-relativo}
\end{equation}

Donde $d^2\vec x_{P/O'}/dt^2$ es la aceleración del cuerpo $P$ con respecto al observador $O'$, la que escribiremos como
\begin{equation}
    \left.\frac{d^2\vec x}{dt^2}\right|_{O'}=\frac{d^2\vec x_{P/O'}}{dt^2}.
    \label{eq:mecanica-aceleracion-relativa-op}
\end{equation}
Introduciendo~\eqref{eq:mecanica-aceleracion-relativa-op} en~\eqref{eq:mecanica-segunda-newton-reemplazo-relativo} se encuentra
\begin{equation}
    m\left(\left.\frac{d^2\vec x}{dt^2}\right|_{O'}+\frac{d^2\vec x_{O'/O}}{dt^2}\right)=\vec F.
    \label{eq:mecanica-segunda-newton-no-inercial-previa}
\end{equation}

En el lado izquierdo de la ecuación~\eqref{eq:mecanica-segunda-newton-no-inercial-previa} aparece la expresión $d^2\vec x_{O'/O}/dt^2$, que es la aceleración del observador $O'$ con respecto al observador $O$. Si pasamos la expresión $d^2\vec x_{O'/O}/dt^2$ al lado izquierdo de la ecuación~\eqref{eq:mecanica-segunda-newton-no-inercial-previa} se obtiene que
\begin{equation}
    m\left.\frac{d^2\vec x}{dt^2}\right|_{O'}=\vec F-m\frac{d^2\vec x_{O'/O}}{dt^2}.
    \label{eq:mecanica-segunda-newton-no-inercial}
\end{equation}

\subsection{Fuerza inercial}

La expresión $-m d^2\vec x_{O'/O}/dt^2$ está en el lado derecho de la ecuación, la que corresponde al lado dinámico por lo que se interpreta como una fuerza y se le denomina como fuerza inercial:
\begin{equation}
    \vec F_{\text{inercial}}=-m\frac{d^2\vec x_{O'/O}}{dt^2},
    \label{eq:mecanica-fuerza-inercial}
\end{equation}
tal que~\eqref{eq:mecanica-segunda-newton-no-inercial} se reescribe como
\begin{equation}
    m\left.\frac{d^2\vec x}{dt^2}\right|_{O'}=\vec F+\vec F_{\text{inercial}}.
    \label{eq:mecanica-segunda-newton-con-fuerza-inercial}
\end{equation}
Para el observador $O'$ el número de fuerzas aumenta,
\begin{equation}
    \left.\vec F\right|_{O'}=\vec F+\vec F_{\text{inercial}}.
    \label{eq:mecanica-fuerza-total-no-inercial}
\end{equation}

La fuerza $\vec F_{\text{inercial}}$ no corresponde, estrictamente hablando, a una fuerza, pues no hay presente un agente externo que genere esta fuerza, sino que es debida solo a los efectos inerciales, o sea a los efectos del espacio absoluto, espacio muy cuestionado.

Como conclusión: los observadores no inerciales observan fuerzas adicionales, que no son estrictamente causales. En lo que sigue consideraremos observadores inerciales, a menos que se diga lo contrario.

\subsection{Trayectoria y parametrización temporal}

El primer paso para hacer uso de la primera ley de Newton es definir un sistema de referencia. El segundo paso es la elección del sistema de coordenadas, que esté en acorde con la geometría que posee el sistema que describe al moverse, y una buena elección resulta muy útil a la hora de simplificar las ecuaciones del movimiento.

Una vez escogido el sistema de coordenado, se procede a definir la posición del cuerpo en el espacio
\begin{equation*}
    \vec x(x^1,x^2,x^3)=\vec x(x^i).
\end{equation*}
Como la partícula se mueve, entonces cambia su posición en el espacio, describiendo una curva continua en el espacio, llamada trayectoria, como se esquematiza en la Figura~\ref{fig:mecanica-trayectoria}.

\begin{figure}[H]
    \centering
    \IfFileExists{figures/08_mecanica/trayectoria_particula.png}{%
        \includegraphics[width=0.45\linewidth]{figures/08_mecanica/trayectoria_particula.png}%
    }{%
        \fbox{\parbox{0.45\linewidth}{\centering Figura pendiente: trayectoria de una partícula.}}%
    }
    \caption{Trayectoria de una partícula en un sistema de coordenadas.}
    \label{fig:mecanica-trayectoria}
\end{figure}

\begin{apuntebloque}{Definición}
La \textbf{trayectoria} que describe una partícula es un conjunto continuo de posiciones.
\end{apuntebloque}

La trayectoria puede ser parametrizada por una magnitud física llamada tiempo:
\begin{equation*}
    \vec x(x^1(t),x^2(t),x^3(t))=\vec x(x^i(t)).
\end{equation*}
Las coordenadas dependen del tiempo, $x^i(t)$. Entonces, la trayectoria es una curva con parámetro $\lambda=t$, que se estudió a principio del curso.

\newpage
\section{Velocidad}

\subsection{Velocidad como vector tangente}

La curva posee un vector tangente, que en mecánica clásica corresponde al llamado vector velocidad:
\begin{equation*}
    \vec v=\frac{d\vec x}{dt},
\end{equation*}
y nos dice hacia donde evoluciona la curva, o sea $\vec v$ apunta en la dirección del movimiento, como se indica en la Figura~\ref{fig:mecanica-vector-velocidad}.

\begin{figure}[H]
    \centering
    \IfFileExists{figures/08_mecanica/vector_velocidad.png}{%
        \includegraphics[width=0.45\linewidth]{figures/08_mecanica/vector_velocidad.png}%
    }{%
        \fbox{\parbox{0.45\linewidth}{\centering Figura pendiente: vector velocidad tangente a la trayectoria.}}%
    }
    \caption{Vector velocidad como vector tangente a la trayectoria.}
    \label{fig:mecanica-vector-velocidad}
\end{figure}

\subsection{Componentes en una base coordenada}

Escogido el sistema de coordenadas se establecen los vectores base coordenadas $\{\vec e_i\}_{i=1}^{3}$, tal que el vector velocidad se escribirá en función de los vectores base como
\begin{equation}
    \vec v=\sum_{i=1}^{3}v^i\vec e_i,
    \label{eq:mecanica-velocidad-base-coordenada}
\end{equation}
donde $v^i\in\mathbb R$ es la componente de la velocidad en la base coordenado $\{\vec e_i\}_{i=1}^{3}$, y puede depender de las coordenadas $x^i$ y por lo tanto del tiempo,
\begin{equation}
    v^i=v^i(x^j(t)).
    \label{eq:mecanica-componentes-velocidad-dependencia}
\end{equation}

La velocidad y la posición están relacionadas en la forma
\begin{align}
    \vec v
    &=\frac{d\vec x}{dt}\bigl(x^j(t)\bigr) \\
    &=\frac{d\vec x}{dt}\bigl(x^1(t),x^2(t),x^3(t)\bigr) \\
    &=\frac{dx^1}{dt}\frac{\partial \vec x}{\partial x^1}
     +\frac{dx^2}{dt}\frac{\partial \vec x}{\partial x^2}
     +\frac{dx^3}{dt}\frac{\partial \vec x}{\partial x^3} \\
    &=\sum_{i=1}^{3}\frac{dx^i}{dt}\frac{\partial \vec x}{\partial x^i}
     =\sum_{i=1}^{3}\frac{dx^i}{dt}\vec e_i .
\end{align}
Por lo tanto,
\begin{equation}
    \vec v=\sum_{i=1}^{3}\frac{dx^i}{dt}\vec e_i.
    \label{eq:mecanica-velocidad-derivadas-coordenadas}
\end{equation}
Comparando~\eqref{eq:mecanica-velocidad-derivadas-coordenadas} con~\eqref{eq:mecanica-velocidad-base-coordenada} se obtiene que
\begin{equation}
    v^i=\frac{dx^i}{dt}.
    \label{eq:mecanica-componentes-velocidad}
\end{equation}

\begin{apuntebloque}{Ejemplo}
\textbf{a) Coordenadas cartesianas.} La velocidad está dada por
\begin{equation*}
    \vec v=\frac{dx}{dt}\hat{\imath}+\frac{dy}{dt}\hat{\jmath}+\frac{dz}{dt}\hat{k},
\end{equation*}
luego
\begin{equation*}
    \frac{dx}{dt}=\dot x,
    \qquad
    \frac{dy}{dt}=\dot y,
    \qquad
    \frac{dz}{dt}=\dot z.
\end{equation*}
\end{apuntebloque}

\begin{apuntebloque}{Ejemplo}
\textbf{b) Coordenadas esféricas.} La velocidad está dada por
\begin{equation*}
    \vec v=\frac{dr}{dt}\vec e_r+\frac{d\phi}{dt}\vec e_\phi+\frac{d\theta}{dt}\vec e_\theta,
\end{equation*}
luego
\begin{equation*}
    \frac{dr}{dt}=\dot r,
    \qquad
    \frac{d\phi}{dt}=\dot\phi,
    \qquad
    \frac{d\theta}{dt}=\dot\theta.
\end{equation*}
\end{apuntebloque}

\subsection{Componentes físicas en una base unitaria}

Es costumbre usar los vectores unitarios coordenados, para que los $v^i$ posean unidades de velocidad. Luego de~\eqref{eq:mecanica-velocidad-base-coordenada} y usando $\vec e_i=\|\vec e_i\|\hat e_i$ se tiene que
\begin{equation*}
    \vec v=\sum_{i=1}^{3}\frac{dx^i}{dt}\|\vec e_i\|\hat e_i
    =\sum_{i=1}^{3}\widetilde v^i\hat e_i,
\end{equation*}
donde se ha efectuado la identificación
\begin{equation}
    \widetilde v^i=\frac{dx^i}{dt}\|\vec e_i\|=\|\vec e_i\|v^i.
    \label{eq:mecanica-componentes-velocidad-unitaria}
\end{equation}

\begin{apuntebloque}{Ejemplo}
\textbf{a) Coordenadas cartesianas.} En coordenadas cartesianas $\|\hat e_i\|=1$, por lo que
\begin{equation*}
    \widetilde v^i=v^i.
\end{equation*}
\end{apuntebloque}

\begin{apuntebloque}{Ejemplo}
\textbf{b) Coordenadas esféricas.} En coordenadas esféricas la velocidad se escribe en términos de los vectores unitarios como
\begin{equation*}
    \vec v=\dot r\hat r+\dot\phi\,r\sin\theta\,\hat\phi+\dot\theta\,r\hat\theta,
\end{equation*}
luego
\begin{equation*}
    \widetilde v^r=\dot r,
    \qquad
    \widetilde v^\phi=r\sin\theta\,\dot\phi,
    \qquad
    \widetilde v^\theta=r\dot\theta.
\end{equation*}
\end{apuntebloque}

\subsection{Rapidez y energía cinética}

La \textbf{rapidez} está dada por el módulo de la velocidad:
\begin{equation}
    v=\|\vec v\|=\sqrt{\sum_{i=1}^{3}\sum_{j=1}^{3}g_{ij}v^i v^j}.
    \label{eq:mecanica-rapidez-general}
\end{equation}

\begin{apuntebloque}{Ejemplo}
\textbf{a) Coordenadas cartesianas.} La rapidez en coordenadas cartesianas es
\begin{align*}
    v
    &=\sqrt{\sum_{i=1}^{3}\sum_{j=1}^{3}\delta_{ij}v^i v^j}
      =\sqrt{\sum_{i=1}^{3}v^i v^i} \\
    &=\sqrt{(v^x)^2+(v^y)^2+(v^z)^2}
      =\sqrt{\dot x^2+\dot y^2+\dot z^2}.
\end{align*}
\end{apuntebloque}

\begin{apuntebloque}{Ejemplo}
\textbf{b) Coordenadas esféricas.} La rapidez en coordenadas esféricas está dada por
\begin{align*}
    v=\|\vec v\|
    &=\sqrt{\sum_{i=1}^{3}\sum_{j=1}^{3}g_{ij}v^i v^j} \\
    &=\sqrt{g_{11}\dot r^2+g_{22}\dot\phi^2+g_{33}\dot\theta^2} \\
    &=\sqrt{\dot r^2+r^2\sin^2\theta\,\dot\phi^2+r^2\dot\theta^2}.
\end{align*}
\end{apuntebloque}

La energía cinética está dada por
\begin{equation*}
    T=\frac{1}{2}mv^2=\frac{1}{2}m\sum_{i=1}^{3}\sum_{j=1}^{3}g_{ij}v^i v^j.
\end{equation*}
Si hay $N$ partículas de masa $m_a$ y coordenadas $x^i_a$, con $a=1,2,\ldots,N$, entonces la energía cinética es dada por
\begin{equation*}
    T=\sum_{a=1}^{N}T_a
    =\frac{1}{2}\sum_{a=1}^{N}m_a
    \sum_{i=1}^{3}\sum_{j=1}^{3}g_{ij}(x^k_a)v^i_a v^j_a,
\end{equation*}
donde las coordenadas curvilíneas son las mismas para todas las partículas.

\newpage
\section{Distancia}

\subsection{Longitud de trayectoria}

\begin{apuntebloque}{Definición}
La \textbf{distancia recorrida} por la partícula al moverse es la longitud de la trayectoria.
\end{apuntebloque}

Como $x^i=x^i(t)$, entonces
\begin{align}
    ds^2
    &=\sum_{i=1}^{3}\sum_{j=1}^{3}g_{ij}dx^i dx^j \\
    &=\sum_{i=1}^{3}\sum_{j=1}^{3}g_{ij}v^i dt\,v^j dt \\
    &=\sum_{i=1}^{3}\sum_{j=1}^{3}g_{ij}v^i v^j dt^2
      =v^2dt^2.
\end{align}
Por lo tanto,
\begin{equation}
    ds=vdt.
    \label{eq:mecanica-ds-vdt}
\end{equation}

Por lo tanto, la distancia recorrida entre $t=t_i$ y $t=t_f$ es
\begin{equation}
    \Delta S=\int_{t_i}^{t_f}vdt.
    \label{eq:mecanica-distancia-integral}
\end{equation}
Notar que
\begin{equation}
    v=\frac{ds}{dt},
    \label{eq:mecanica-rapidez-dsdt}
\end{equation}
la rapidez corresponde al cambio de la distancia en el tiempo.

\subsection{Lanzamiento vertical}

\begin{apuntebloque}{Ejemplo}
\textbf{a)} En un lanzamiento vertical un cuerpo es lanzado hacia arriba en una línea recta. La ecuación del movimiento es dada por
\begin{equation*}
\left.
\begin{aligned}
    x(t)&=0,\\
    y(t)&=v_0t-\frac{1}{2}gt^2,\\
    z(t)&=0,
\end{aligned}
\right\}
\qquad\Longrightarrow\qquad
\left\{
\begin{aligned}
    v_x(t)&=0,\\
    v_y(t)&=v_0-gt,\\
    v_z(t)&=0,
\end{aligned}
\right.
\end{equation*}
con posición inicial cero y $v_0$ es la velocidad inicial.
\end{apuntebloque}

El tiempo en alcanzar la altura máxima es cuando la velocidad es cero;
\begin{equation*}
    v_y(t)=v_0-gt=0,
\end{equation*}
por lo tanto
\begin{equation}
    t_H=\frac{v_0}{g}.
    \label{eq:mecanica-tiempo-altura-maxima-vertical}
\end{equation}
Y la altura máxima es dada por
\begin{align}
    H=y(t_H)&=v_0t_H-\frac{g}{2}t_H^2
    =v_0\cdot\frac{v_0}{g}-\frac{g}{2}\cdot\frac{v_0^2}{g^2},
\end{align}
de donde
\begin{equation}
    H=\frac{v_0^2}{2g}.
    \label{eq:mecanica-altura-maxima-vertical}
\end{equation}

Por lo tanto, la distancia que sube la partícula es
\begin{equation}
    (\Delta S)_{\text{subida}}=H=\frac{v_0^2}{2g},
    \label{eq:mecanica-distancia-subida-vertical}
\end{equation}
y por como consecuencia la distancia que recorre de bajada también es $H$,
\begin{equation}
    (\Delta S)_{\text{bajada}}=H=\frac{v_0^2}{2g}.
    \label{eq:mecanica-distancia-bajada-vertical}
\end{equation}
Así, la distancia total recorrida por la partícula es
\begin{equation*}
    \Delta S=(\Delta S)_{\text{subida}}+(\Delta S)_{\text{bajada}}=H+H=2H,
\end{equation*}
por lo tanto
\begin{equation}
    \Delta S=\frac{v_0^2}{g}.
    \label{eq:mecanica-distancia-total-vertical}
\end{equation}

Ahora calculemos la distancia usando la ecuación~\eqref{eq:mecanica-distancia-integral}, para ello calculemos el tiempo de vuelo:
\begin{equation*}
    y(t)=v_0t-\frac{1}{2}gt^2=0.
\end{equation*}
Se obtienen dos tiempos
\begin{equation}
    t_i=0,
    \qquad
    t_v=\frac{2v_0}{g}=2t_H.
    \label{eq:mecanica-tiempo-vuelo-vertical}
\end{equation}

La distancia recorrida es
\begin{align}
    \Delta S
    &=\int_{0}^{t_v}vdt
      =\int_{0}^{t_v}\sqrt{v_x^2+v_y^2+v_z^2}\,dt \\
    &=\int_{0}^{t_v}\sqrt{0^2+v_y^2+0^2}\,dt
      =\int_{0}^{t_v}v_y dt \\
    &=\int_{0}^{t_v}(v_0-gt)dt
      =\left(v_0t-\frac{1}{2}gt^2\right)\Big|_{0}^{t_v} \\
    &=y(t_v)-y(0)=0+0.
\end{align}
Por lo tanto,
\begin{equation}
    \Delta S=0.
    \label{eq:mecanica-distancia-cancelacion-vertical}
\end{equation}

Claramente hay una discrepancia entre los resultados~\eqref{eq:mecanica-distancia-total-vertical} y~\eqref{eq:mecanica-distancia-cancelacion-vertical}. La razón de esta discrepancia es debida a que la rapidez debe satisfacer $v\geq 0$, al igual que la distancia, y para este problema se tiene que $v=v_y$, y grafiquemos $v_y$. La interpretación geométrica de esta cancelación se indica en la Figura~\ref{fig:mecanica-lanzamiento-vertical-areas}.

\begin{figure}[H]
    \centering
    \IfFileExists{figures/08_mecanica/lanzamiento_vertical_areas.png}{%
        \includegraphics[width=0.55\linewidth]{figures/08_mecanica/lanzamiento_vertical_areas.png}%
    }{%
        \fbox{\parbox{0.55\linewidth}{\centering Figura pendiente: gráfico $v_y$ vs. $t$ y áreas $A_1$, $A_2$.}}%
    }
    \caption{Gráfico de $v_y$ contra $t$ en el lanzamiento vertical.}
    \label{fig:mecanica-lanzamiento-vertical-areas}
\end{figure}

Por lo tanto, de la gráfica $v_y$ vs. $t$ se concluye que $v_y\leq 0$ para $t_H\leq t\leq t_v$, entonces la integral~\eqref{eq:mecanica-distancia-integral}, con $v=v_y$, será la suma de un área $A_1$ positiva y un área $A_2$ negativa, suma que será igual a cero, puesto que $A_2=-A_1$.

Lo anteriormente explicado implica que la rapidez es negativa en el intervalo de tiempo $t_H\leq t\leq t_v$, lo que es una contradicción, puesto que la rapidez, por definición, es positiva. La definición de la rapidez involucra la raíz cuadrada positiva, $+\sqrt{\cdots}$, y no la raíz negativa, $-\sqrt{\cdots}$, luego se tiene que
\begin{align}
    v
    &=\sqrt{(v_x)^2+(v_y)^2+(v_z)^2} \\
    &=\sqrt{(0)^2+(v_y)^2+(0)^2}
      =\sqrt{(v_y)^2}.
\end{align}
Por lo tanto,
\begin{equation}
    v(t)=|v_y(t)|.
    \label{eq:mecanica-rapidez-valor-absoluto}
\end{equation}

La ecuación~\eqref{eq:mecanica-rapidez-valor-absoluto} se traduce a
\begin{equation}
    v(t)=
    \begin{cases}
        v_y, & 0\leq t\leq t_H,\\
        -v_y, & t_H\leq t\leq t_v,
    \end{cases}
    =
    \begin{cases}
        v_0-gt, & 0\leq t\leq t_H,\\
        gt-v_0, & t_H\leq t\leq t_v.
    \end{cases}
    \label{eq:mecanica-rapidez-vertical-por-tramos}
\end{equation}

Por lo tanto, la distancia recorrida por la partícula es
\begin{align*}
    \Delta S
    &=\int_{t_i}^{t_f}vdt
      =\int_{0}^{t_v}|v_y|dt
      =\int_{0}^{t_H}v_y dt+\int_{t_H}^{t_v}(-v_y)dt \\
    &=\int_{0}^{t_H}(v_0-gt)dt+\int_{t_H}^{t_v}(gt-v_0)dt \\
    &=\left(v_0t-\frac{1}{2}gt^2\right)\Big|_{0}^{t_H}
      +\left(\frac{1}{2}gt^2-v_0t\right)\Big|_{t_H}^{t_v} \\
    &=y(t)\Big|_{0}^{t_H}-y(t)\Big|_{t_H}^{t_v} \\
    &=y(t_H)-y(0)-\bigl(y(t_v)-y(t_H)\bigr) \\
    &=2y(t_H)-y(0)-y(t_v)
      =2y(t_H)=\frac{v_0^2}{g}.
\end{align*}
Se recupera~\eqref{eq:mecanica-distancia-total-vertical}.

\subsection{Movimiento parabólico}

\begin{apuntebloque}{Ejemplo}
\textbf{b)} En un movimiento parabólico, un cuerpo es lanzado desde el origen, con rapidez inicial $v_0$ y ángulo de elevación $\theta$. La ecuación del movimiento es dada por
\begin{equation}
\left.
\begin{aligned}
    x(t)&=v_0\cos\theta\,t,\\
    y(t)&=v_0\sin\theta\,t-\frac{1}{2}gt^2,\\
    z(t)&=0,
\end{aligned}
\right\}
\qquad\Longrightarrow\qquad
\left\{
\begin{aligned}
    v_x(t)&=v_0\cos\theta,\\
    v_y(t)&=v_0\sin\theta-gt,\\
    v_z(t)&=0,
\end{aligned}
\right.
\label{eq:mecanica-movimiento-parabolico}
\end{equation}
con posición inicial cero y donde $v_0$ es la rapidez inicial.
\end{apuntebloque}

El tiempo en alcanzar la altura máxima es cuando la componente de la velocidad $v_y$ es cero;
\begin{equation*}
    v_y=v_0\sin\theta-gt=0,
\end{equation*}
por lo tanto
\begin{equation}
    t_H=\frac{v_0\sin\theta}{g}.
    \label{eq:mecanica-tiempo-altura-maxima-parabolico}
\end{equation}
Y la altura máxima es dada por
\begin{align*}
    H=y(t_H)
    &=v_0\sin\theta\,t_H-\frac{1}{2}gt_H^2 \\
    &=v_0\sin\theta\cdot\frac{v_0\sin\theta}{g}
      -\frac{1}{2}g\left(\frac{v_0\sin\theta}{g}\right)^2.
\end{align*}
Luego,
\begin{equation}
    H=\frac{v_0^2\sin^2\theta}{2g}.
    \label{eq:mecanica-altura-maxima-parabolico}
\end{equation}

Calculemos el tiempo de vuelo que nos permitirá calcular la distancia recorrida por la partícula:
\begin{align*}
    y(t)&=0=v_0\sin\theta\,t-\frac{1}{2}gt^2
      =\left(v_0\sin\theta-\frac{1}{2}gt\right)t.
\end{align*}
Se obtienen dos tiempos,
\begin{equation}
    t_i=0,
    \qquad
    t_v=\frac{2v_0\sin\theta}{g}=2t_H.
    \label{eq:mecanica-tiempo-vuelo-parabolico}
\end{equation}
El alcance máximo es
\begin{equation}
    R=x(t_v)=v_0\cos\theta\,t_v
      =v_0\cos\theta\cdot\frac{2v_0\sin\theta}{g}
      =\frac{v_0^2\sin(2\theta)}{g}.
    \label{eq:mecanica-alcance-parabolico}
\end{equation}

Ahora calculemos la distancia recorrida por la partícula, esto es la longitud de la parábola:
\begin{align*}
    \Delta S
    &=\int_{t_i}^{t_f}vdt
      =\int_{0}^{t_v}\sqrt{v_x^2+v_y^2+v_z^2}\,dt \\
    &=\int_{0}^{t_v}\sqrt{(v_0\cos\theta)^2+(v_0\sin\theta-gt)^2+(0)^2}\,dt.
\end{align*}
Por lo tanto,
\begin{equation}
    \Delta S=\frac{v_0^2}{g}
    \left(\sin\theta+(\cos\theta)^2\ln\left|\frac{1+\sin\theta}{\cos\theta}\right|\right).
    \label{eq:mecanica-longitud-parabolica}
\end{equation}

La distancia depende de la rapidez inicial y del ángulo de elevación. Definamos la función
\begin{equation}
    f(\theta)=\sin\theta+(\cos\theta)^2\ln\left|\frac{1+\sin\theta}{\cos\theta}\right|.
    \label{eq:mecanica-funcion-f-theta}
\end{equation}
Luego reemplazando en~\eqref{eq:mecanica-longitud-parabolica} se obtiene
\begin{equation}
    \Delta S(v_0,\theta)=\frac{v_0^2}{g}f(\theta).
    \label{eq:mecanica-distancia-f-theta}
\end{equation}

Analicemos la expresión~\eqref{eq:mecanica-distancia-f-theta}: si dejamos fijo el ángulo de elevación $\theta$ y variamos la rapidez inicial $v_0$, la gráfica $\Delta S$ vs. $v_0$ corresponde a una parábola creciente desde cero. A medida que $v_0$ aumenta la distancia aumenta con la cuadratura de la rapidez, y la distancia crece sin cota si $v_0$ crece sin cota, como se representa en la Figura~\ref{fig:mecanica-distancia-v0}.

\begin{figure}[H]
    \centering
    \IfFileExists{figures/08_mecanica/distancia_vs_v0.png}{%
        \includegraphics[width=0.45\linewidth]{figures/08_mecanica/distancia_vs_v0.png}%
    }{%
        \fbox{\parbox{0.45\linewidth}{\centering Figura pendiente: gráfico $\Delta S$ vs. $v_0$.}}%
    }
    \caption{Comportamiento de $\Delta S$ al variar la rapidez inicial $v_0$.}
    \label{fig:mecanica-distancia-v0}
\end{figure}

Ahora dejamos fijo $v_0$ y variamos el ángulo de elevación $\theta$. Para ver como cambia la distancia con el ángulo graficamos la función $f(\theta)$ y se obtiene la Figura~\ref{fig:mecanica-f-theta}.

\begin{figure}[H]
    \centering
    \IfFileExists{figures/08_mecanica/funcion_f_theta.png}{%
        \includegraphics[width=0.55\linewidth]{figures/08_mecanica/funcion_f_theta.png}%
    }{%
        \fbox{\parbox{0.55\linewidth}{\centering Figura pendiente: gráfico $f(\theta)$ vs. $\theta$.}}%
    }
    \caption{Gráfico $f(\theta)$ vs. $\theta$.}
    \label{fig:mecanica-f-theta}
\end{figure}

Del gráfico $f(\theta)$ vs. $\theta$ se observa que la función $f(\theta)$ crece desde cero, para $\theta=0$, a un valor máximo y luego decrece hasta que $\theta$ toma el valor $\pi/2$. Por lo tanto, la función $f(\theta)$ posee un valor máximo (que no es $\pi/4$ donde se obtiene el alcance máximo) para algún valor de $\theta$ y como consecuencia para ese valor la longitud de la parábola es máxima.

Para encontrar el valor crítico de $\theta$ debemos derivar $f(\theta)$ e igualarla a cero, y se obtiene
\begin{equation*}
    \frac{df}{d\theta}=\cos\theta\left(1-\sin\theta\ln\left|\frac{1+\sin\theta}{\cos\theta}\right|\right)=0.
\end{equation*}
Tenemos que resolver la ecuación trascendental
\begin{equation}
    1-\sin\theta\ln\left|\frac{1+\sin\theta}{\cos\theta}\right|=0.
    \label{eq:mecanica-ecuacion-critica-f-theta}
\end{equation}

La representación gráfica nos indica que el máximo de $f(\theta)$ se encuentra entre $50^{\circ}$ y $60^{\circ}$. Se calcula la raíz de la ecuación trascendente por el procedimiento del punto medio. El valor que se obtiene es $\theta_m=56.46^{\circ}=0.31\pi\,\mathrm{rad}$.

Entonces el valor máximo de la la función $f(\theta)$ es
\begin{equation*}
    f_{\max}=f(\theta_m)=1.199515690,
\end{equation*}
tal que la parábola que posee la longitud máxima es de valor
\begin{equation*}
    \Delta S_{\max}=1.199515690\cdot\frac{v_0^2}{g}.
\end{equation*}

Para $\theta=90^{\circ}$ la función logaritmo natural de la función $f(\theta)$ se indetermina, puesto que $1/\cos\theta$ se indetermina. Sin embargo, la función logaritmo natural está multiplicado por $(\cos\theta)^2$ el que toma el valor cero, y por lo tanto la función $f(\theta)$ es suave en ese valor. De la gráfica $f(\theta)$ vs. $\theta$ se obtiene el valor
\begin{equation*}
    f(90^{\circ})=1.
\end{equation*}
Tal que
\begin{equation*}
    \Delta S(v_0,90^{\circ})=\frac{v_0^2}{g},
\end{equation*}
que corresponde al lanzamiento vertical.

¿Qué pasa con el caso $\theta=0$?

\newpage
\section{Aceleración}

\subsection{Derivada temporal de la velocidad}

La \textbf{aceleración} corresponde al cambio de la velocidad en función del tiempo, así
\begin{align*}
    \vec a
    &=\frac{d\vec v}{dt}
      =\frac{d}{dt}\left(\sum_{i=1}^{3}v^i\vec e_i\right) \\
    &=\sum_{i=1}^{3}\left(\frac{dv^i}{dt}\vec e_i+v^i\frac{d\vec e_i}{dt}\right).
\end{align*}

La base coordenada $\{\vec e_i\}_{i=1}^{3}$ depende implícitamente del tiempo, a través de las coordenadas $x^i(t)$, tal que
\begin{equation*}
    \vec e_i=\vec e_i(x^i(t)),
\end{equation*}
y no depende explícitamente del tiempo (sistema de referencia inercial), luego
\begin{align*}
    \frac{d\vec e_i}{dt}
    &=\frac{dx^1}{dt}\frac{\partial \vec e_i}{\partial x^1}
     +\frac{dx^2}{dt}\frac{\partial \vec e_i}{\partial x^2}
     +\frac{dx^3}{dt}\frac{\partial \vec e_i}{\partial x^3} \\
    &=\sum_{j=1}^{3}\frac{dx^j}{dt}\frac{\partial \vec e_i}{\partial x^j} \\
    &=\sum_{j=1}^{3}\frac{dx^j}{dt}\sum_{k=1}^{3}\Gamma^{k}_{ij}\vec e_k \\
    &=\sum_{j=1}^{3}\sum_{k=1}^{3}\Gamma^{k}_{ij}v^j\vec e_k,
\end{align*}
se usó~\eqref{eq:mecanica-componentes-velocidad}.

Entonces
\begin{align*}
    \vec a
    &=\sum_{i=1}^{3}\left(\frac{dv^i}{dt}\vec e_i+v^i\sum_{j=1}^{3}\sum_{k=1}^{3}\Gamma^{k}_{ij}v^j\vec e_k\right) \\
    &=\sum_{i=1}^{3}\frac{dv^i}{dt}\vec e_i
      +\sum_{i=1}^{3}\sum_{j=1}^{3}\sum_{k=1}^{3}\Gamma^{k}_{ij}v^iv^j\vec e_k \\
    &=\sum_{i=1}^{3}\frac{dv^i}{dt}\vec e_i
      +\sum_{k=1}^{3}\sum_{j=1}^{3}\sum_{i=1}^{3}\Gamma^{i}_{kj}v^kv^j\vec e_i.
\end{align*}
Por lo tanto,
\begin{equation}
    \vec a=\sum_{i=1}^{3}\left(\frac{dv^i}{dt}+\sum_{j=1}^{3}\sum_{k=1}^{3}\Gamma^{i}_{jk}v^jv^k\right)\vec e_i.
    \label{eq:mecanica-aceleracion-general}
\end{equation}
\subsection{Componentes de la aceleración}

Como la aceleración es un vector, entonces
\begin{equation}
    \vec a=\sum_{i=1}^{3}a^i\vec e_i.
    \label{eq:mecanica-aceleracion-componentes}
\end{equation}
Por lo tanto, se concluye que
\begin{equation}
    a^i=\frac{dv^i}{dt}+\sum_{j=1}^{3}\sum_{k=1}^{3}\Gamma^{i}_{jk}v^jv^k.
    \label{eq:mecanica-componentes-aceleracion-general}
\end{equation}

La componente $dv^i/dt$ mide como cambia la magnitud de la velocidad. La componente $\sum_{j=1}^{3}\sum_{k=1}^{3}\Gamma^{i}_{jk}v^jv^k$ mide como cambia la dirección el vector velocidad.

\begin{apuntebloque}{Ejemplo}
\textbf{a) Coordenadas cartesianas.} En las coordenadas cartesianas tenemos que $\Gamma^{i}_{jk}=0$, luego
\begin{equation*}
    a^x=\frac{dv^x}{dt}=\frac{d^2x}{dt^2},
    \qquad
    a^y=\frac{dv^y}{dt}=\frac{d^2y}{dt^2},
    \qquad
    a^z=\frac{dv^z}{dt}=\frac{d^2z}{dt^2}.
\end{equation*}
\end{apuntebloque}

\begin{apuntebloque}{Ejemplo}
\textbf{b) Coordenadas esféricas.} En coordenadas esféricas tenemos que la aceleración es
\begin{align*}
    a^r&=\ddot r-r\sin^2\theta\,\dot\phi^2-r\dot\theta^2,\\
    a^\phi&=\ddot\phi+\frac{2}{r}\dot r\dot\phi+2\cot\theta\,\dot\phi\dot\theta,\\
    a^\theta&=\ddot\theta-\frac{1}{2}\sin(2\theta)\dot\phi^2+\frac{2}{r}\dot r\dot\theta.
\end{align*}
\end{apuntebloque}

\subsection{Componentes físicas en una base unitaria}

Es costumbre usar los vectores unitarios coordenados, para que los $a^i$ posean unidades de aceleración. Luego de~\eqref{eq:mecanica-velocidad-base-coordenada} y usando $\vec e_i=\|\vec e_i\|\hat e_i$ se tiene que
\begin{equation*}
    \vec a=\sum_{i=1}^{3}a^i\|\vec e_i\|\hat e_i=\sum_{i=1}^{3}\widetilde a^i\hat e_i,
\end{equation*}
donde se efectuado la identificación
\begin{equation}
    \widetilde a^i=\|\vec e_i\|a^i
    =\left(\frac{dv^i}{dt}+\sum_{j=1}^{3}\sum_{k=1}^{3}\Gamma^{i}_{jk}v^jv^k\right)\|\vec e_i\|.
    \label{eq:mecanica-componentes-aceleracion-unitaria}
\end{equation}

\begin{apuntebloque}{Ejemplo}
\textbf{a) Coordenadas cartesianas.} En coordenadas cartesianas $\|\hat e_i\|=1$, por lo que
\begin{equation*}
    \widetilde a^i=a^i.
\end{equation*}
\end{apuntebloque}

\begin{apuntebloque}{Ejemplo}
\textbf{b) Coordenadas esféricas.} En coordenadas esféricas la aceleración se escribe en términos de los vectores unitarios como
\begin{align*}
    \vec a
    &=\left(\ddot r-r\sin^2\theta\,\dot\phi^2-r\dot\theta^2\right)\hat r \\
    &\quad+
    \left(\ddot\phi+\frac{2}{r}\dot r\dot\phi+2\cot\theta\,\dot\phi\dot\theta\right)r\sin\theta\,\hat\phi \\
    &\quad+
    \left(\ddot\theta-\frac{1}{2}\sin(2\theta)\dot\phi^2+\frac{2}{r}\dot r\dot\theta\right)r\hat\theta.
\end{align*}
Luego
\begin{align*}
    \widetilde a^r&=\ddot r-r\sin^2\theta\,\dot\phi^2-r\dot\theta^2,\\
    \widetilde a^\phi&=\left(\ddot\phi+\frac{2}{r}\dot r\dot\phi+2\cot\theta\,\dot\phi\dot\theta\right)r\sin\theta,\\
    \widetilde a^\theta&=\left(\ddot\theta-\frac{1}{2}\sin(2\theta)\dot\phi^2+\frac{2}{r}\dot r\dot\theta\right)r.
\end{align*}
\end{apuntebloque}

\subsection{Ecuaciones de movimiento en coordenadas generales}

La segunda ley de Newton toma la forma
\begin{equation*}
    m\sum_{i=1}^{3}\left(\frac{dv^i}{dt}+\sum_{j=1}^{3}\sum_{k=1}^{3}\Gamma^{i}_{jk}v^jv^k\right)\vec e_i
    =\sum_{i=1}^{3}F^i\vec e_i.
\end{equation*}
Y puesto que las bases coordenadas son linealmente independientes, entonces
\begin{equation}
    m\left(\frac{dv^i}{dt}+\sum_{j=1}^{3}\sum_{k=1}^{3}\Gamma^{i}_{jk}v^jv^k\right)=F^i.
    \label{eq:mecanica-segunda-newton-curvilinea}
\end{equation}

Conocida la fuerza, caso general, en función de la posición, velocidad y del tiempo,
\begin{equation}
    F^i=F^i(x^j,v^k,t),
    \label{eq:mecanica-fuerza-general}
\end{equation}
entonces la segunda ley de Newton corresponde a un sistema de ecuaciones diferenciales ordinarias de segundo orden, tal vez no lineales, las que pueden ser resueltas con apropiadas condiciones iniciales
\begin{equation}
    x^i_0=x^i(0),
    \qquad
    v^i_0=v^i(0).
    \label{eq:mecanica-condiciones-iniciales}
\end{equation}
Conocida la posición en función del tiempo, podemos conocer las magnitudes físicas en función del tiempo.

\section{Energía cinética, potencia y trabajo}

\subsection{Potencia y energía cinética}

Si contraemos la segunda ley de Newton con el vector velocidad tenemos que
\begin{equation}
    m\vec a\cdot \vec v=\vec F\cdot\vec v.
    \label{eq:mecanica-contraccion-velocidad}
\end{equation}
De este producto se obtiene la siguiente proposición.

\begin{apuntebloque}{Proposición}
El producto $m\vec a\cdot\vec v$ toma la forma
\begin{equation}
    m\vec a\cdot\vec v=\frac{dT}{dt},
    \label{eq:mecanica-derivada-energia-cinetica}
\end{equation}
donde $T$ es la energía cinética.
\end{apuntebloque}

La expresión~\eqref{eq:mecanica-contraccion-velocidad} con ayuda de~\eqref{eq:mecanica-derivada-energia-cinetica} se reescribe como
\begin{equation}
    \vec F\cdot\vec v=\frac{dT}{dt}.
    \label{eq:mecanica-potencia-teorema-fuerzas-vivas}
\end{equation}
La cantidad $\vec F\cdot\vec v$ corresponde a la llamada potencia:
\begin{equation}
    P=\vec F\cdot\vec v.
    \label{eq:mecanica-definicion-potencia}
\end{equation}
Entonces, la ecuación~\eqref{eq:mecanica-potencia-teorema-fuerzas-vivas} se reescribe como
\begin{equation}
    P=\frac{dT}{dt}.
    \label{eq:mecanica-potencia-cambio-energia-cinetica}
\end{equation}

\subsection{Trabajo mecánico}

\begin{apuntebloque}{Definición}
Se llama \textbf{trabajo mecánico} de una fuerza $\vec F$ a la magnitud física denotado $W$ y definida como
\begin{equation}
    W=\int_C \vec F\cdot d\vec x,
    \label{eq:mecanica-trabajo-mecanico}
\end{equation}
donde $C$ es la trayectoria.
\end{apuntebloque}

El trabajo mecánico mide la capacidad que tiene una fuerza para mover un cuerpo. El trabajo mecánico es la integral de camino desde la posición $A$ hasta la posición $B$. En general, el trabajo depende del camino, o sea, caminos distintos informarán trabajos distintos, como se esquematiza en la Figura~\ref{fig:mecanica-trabajo-camino}.

\begin{figure}[H]
    \centering
    \IfFileExists{figures/08_mecanica/trabajo_mecanico_camino.png}{%
        \includegraphics[width=0.5\linewidth]{figures/08_mecanica/trabajo_mecanico_camino.png}%
    }{%
        \fbox{\parbox{0.5\linewidth}{\centering Figura pendiente: trabajo mecánico como integral de camino.}}%
    }
    \caption{Trabajo mecánico como integral de camino entre dos posiciones.}
    \label{fig:mecanica-trabajo-camino}
\end{figure}

\begin{apuntebloque}{Teorema}
\textbf{Teorema de las fuerzas vivas.} El trabajo mecánico realizado por la fuerza neta es igual al cambio de energía cinética.
\end{apuntebloque}

\emph{Demostración.} Se tiene que
\begin{align*}
    W
    &=\int_A^B \vec F\cdot d\vec x \\
    &=\int_{t_i}^{t_f}m\vec a\cdot\frac{d\vec x}{dt}dt \\
    &=\int_{t_i}^{t_f}m\frac{d\vec v}{dt}\cdot\vec v\,dt \\
    &=\int_{t_i}^{t_f}m\frac{d}{dt}\left(\frac{1}{2}\vec v\cdot\vec v\right)dt \\
    &=\int_{t_i}^{t_f}\frac{d}{dt}\left(\frac{m}{2}v^2\right)dt \\
    &=\int_{t_i}^{t_f}\frac{dT}{dt}dt.
\end{align*}
Por lo tanto,
\begin{equation}
    W=T_f-T_i.
    \label{eq:mecanica-trabajo-cambio-energia-cinetica}
\end{equation}
Q. E. D.

\section{Fuerzas conservativas y energía potencial}

\begin{apuntebloque}{Definición}
Si un campo de fuerza es tal que el trabajo mecánico es el mismo para todos las trayectorias físicamente posibles entre los puntos $A$ y $B$, entonces se dice que la fuerza es conservativa.
\end{apuntebloque}

Consideremos dos trayectorias distintas entre los punto $A$ y $B$. Si la fuerza es conservativa, entonces
\begin{align*}
    \int_{C_1}\vec F\cdot d\vec x
    &=\int_{C_2}\vec F\cdot d\vec x,\\
    \int_{C_1}\vec F\cdot d\vec x
    -\int_{C_2}\vec F\cdot d\vec x&=0,\\
    \int_{C_1}\vec F\cdot d\vec x
    +\int_{C_2}\vec F\cdot d(-\vec x)&=0.
\end{align*}
La cantidad $-d\vec x$ corresponde a recorrer la trayectoria $C_2$ en sentido opuesto. Por lo tanto, si $C_3=-C_2$, esto es equivalente a recorrer una trayectoria cerrada, como se indica en la Figura~\ref{fig:mecanica-camino-cerrado}.

\begin{figure}[H]
    \centering
    \IfFileExists{figures/08_mecanica/caminos_conservativos.png}{%
        \includegraphics[width=0.45\linewidth]{figures/08_mecanica/caminos_conservativos.png}%
    }{%
        \fbox{\parbox{0.45\linewidth}{\centering Figura pendiente: caminos $C_1$, $C_2$ y camino cerrado.}}%
    }
    \caption{Trayectorias distintas entre $A$ y $B$ y construcción de una trayectoria cerrada.}
    \label{fig:mecanica-camino-cerrado}
\end{figure}

Con $C_3=-C_2$,
\begin{align*}
    \int_{C_1}\vec F\cdot d\vec x+
    \int_{-C_2}\vec F\cdot d\vec x&=0,\\
    \int_{C_1}\vec F\cdot d\vec x+
    \int_{C_3}\vec F\cdot d\vec x&=0.
\end{align*}
Esto es equivalente a recorrer una trayectoria cerrada,
\begin{equation*}
    \oint \vec F\cdot d\vec x=0.
\end{equation*}

Si el trabajo de un campo de fuerza a lo largo de un camino cerrado es cero, entonces la fuerza es conservativa, esto es
\begin{equation*}
    \oint \vec F\cdot d\vec x=0.
\end{equation*}

\begin{apuntebloque}{Teorema}
Si $\vec F$ es un campo fuerza conservativo, entonces existe un campo escalar $U(x^i)$ tal que
\begin{equation*}
    \vec F=-\vec\nabla U,
\end{equation*}
donde $U(x^i)$ es la energía potencial y las superficies $U(x^i)=\mathrm{constante}$ se llaman superficies equipotenciales.
\end{apuntebloque}

\emph{Demostración.} Si $\vec F$ es conservativa, entonces
\begin{equation*}
    \oint \vec F\cdot d\vec x=0.
\end{equation*}
Del teorema de Stokes,
\begin{equation*}
    \oint \vec F\cdot d\vec x=\int_S \vec\nabla\times\vec F\cdot d\vec S=0.
\end{equation*}
Puesto que la superficie $S$ es arbitrario, entonces
\begin{equation*}
    \vec\nabla\times\vec F=\vec 0.
\end{equation*}
Del conocido resultado del cálculo diferencial $\vec\nabla\times\vec\nabla\phi=0$, tenemos
\begin{equation*}
    \vec F=-\vec\nabla U.
\end{equation*}
Q. E. D.

En un sistema coordenado general se tiene que
\begin{equation*}
    \vec F=-\sum_{i=1}^{3}\vec E^{i}\frac{\partial U}{\partial x^i}.
\end{equation*}

Para conocer $U(x^i)$ debemos integrar la expresión
\begin{equation*}
    \vec F=-\vec\nabla U
\end{equation*}
a lo largo de una trayectoria, puesto que el potencial nace de una integral de camino, pero, como es conservativa puede ser cualquier trayectoria que conecte los puntos iniciales y finales.

Como queremos conocer $U(x^i)$, el punto final debe ser $x^i$ y el punto inicial debe ser un punto de referencia $x^i_r$, un punto donde se conozca o se defina el valor de la energía potencial. Físicamente, todo lo anterior corresponde a mover una partícula de prueba desde el punto $x^i_r$ hasta el punto $x^i$.

Por lo tanto,
\begin{align*}
    \int_{x^i_r}^{x^i}\vec\nabla U\cdot d\vec x
    &=-\int_{x^i_r}^{x^i}\vec F\cdot d\vec x,\\
    \int_{x^i_r}^{x^i}dU
    &=-\int_{x^i_r}^{x^i}\vec F\cdot d\vec x,\\
    U(x^i)-U(x^i_r)
    &=-\int_{x^i_r}^{x^i}\vec F\cdot d\vec x.
\end{align*}
Por lo tanto,
\begin{equation*}
    U(x^i)=U(x^i_r)-\int_{x^i_r}^{x^i}\vec F\cdot d\vec x.
\end{equation*}

La cantidad $U(x^i_r)$ es el valor que toma la energía potencial en el punto de referencia $x^i_r$. La constante $U(x^i_r)$ aparece porque las cantidades $U$ y $U+\text{constante}$ generan la misma fuerza, por lo tanto la energía potencial puede tomar el valor que queremos con sólo sumarle una constante, por lo tanto el potencial es ambiguo.

Para evitar la ambigüedad anterior, el punto de referencia $x^i_r$ es tal que el potencial sea cero en ese punto, $U(x^i_r)=0$, así
\begin{equation*}
    U(x^i)=-\int_{x^i_r}^{x^i}\vec F\cdot d\vec x.
\end{equation*}

\begin{apuntebloque}{Ejemplo}
Sea una fuerza inversamente proporcional al cuadrado de la distancia,
\begin{equation*}
    \vec F=-\frac{k}{r^2}\hat r.
\end{equation*}
Es fácil mostrar que satisface $\vec\nabla\times\vec F=\vec 0$, entonces debemos escoger el punto de referencia de la energía potencial. Es claro que para $r\to\infty$ la fuerza decae a cero, entonces se escoge el punto de referencia el infinito y se define
\begin{equation*}
    \lim_{r\to\infty}U\to 0.
\end{equation*}
Así,
\begin{align*}
    U
    &=-\int_{\infty}^{r}-\frac{k}{r^2}\hat r\cdot d\vec x
      =\int_{\infty}^{r}\frac{k}{r^2}dr \\
    &=\left.-\frac{k}{r}\right|_{\infty}^{r}.
\end{align*}
Luego,
\begin{equation*}
    U(r)=-\frac{k}{r}.
\end{equation*}
\end{apuntebloque}

Volviendo a la ecuación
\begin{equation*}
    \frac{dT}{dt}=\vec F\cdot\vec v,
\end{equation*}
supongamos que la fuerza es conservativa, entonces
\begin{align*}
    \frac{dT}{dt}
    &=-\sum_{i=1}^{3}\vec E^i\frac{\partial U}{\partial x^i}\cdot
      \sum_{j=1}^{3}v^j\vec e_j \\
    &=-\sum_{i=1}^{3}\sum_{j=1}^{3}\frac{dx^j}{dt}
      \frac{\partial U}{\partial x^i}\vec E^i\cdot\vec e_j \\
    &=-\sum_{i=1}^{3}\sum_{j=1}^{3}\frac{dx^j}{dt}
      \frac{\partial U}{\partial x^i}\delta^i_j \\
    &=-\sum_{i=1}^{3}\frac{dx^i}{dt}\frac{\partial U}{\partial x^i}
      =-\frac{dU}{dt}.
\end{align*}
Por lo tanto,
\begin{align*}
    \frac{dT}{dt}+\frac{dU}{dt}&=0,\\
    \frac{d}{dt}(T+U)&=0,\\
    E=T+U&=\text{constante},
\end{align*}
donde $E$ es la energía del sistema y es una constante del movimiento si la fuerza es conservativa.

\section{Constantes del movimiento}

\begin{apuntebloque}{Definición}
Una cantidad que depende de las coordenadas, velocidad y del tiempo,
\begin{equation*}
    C=C(x^i,v^j,t),
\end{equation*}
se dice que es una constante del movimiento de un sistema si tiene el mismo valor durante toda la evolución del sistema.
\end{apuntebloque}

La definición implica que la constante del movimiento satisface
\begin{align*}
    C(x^i(t_i),v^j(t_i),t_i)
    &=C(x^i(t_1),v^j(t_1),t_1) \\
    &=C(x^i(t_2),v^j(t_2),t_2) \\
    &\quad\vdots \\
    &=C(x^i(t_f),v^j(t_f),t_f).
\end{align*}
Por lo tanto, la constante del movimiento puede ser evaluada en cualquier instante del movimiento. Se acostumbra a evaluarla en el tiempo inicial.

\begin{apuntebloque}{Proposición}
Si $C(x^i,v^j,t)$ es una constante del movimiento, entonces
\begin{equation*}
    \frac{dC}{dt}=0.
\end{equation*}
\end{apuntebloque}

\emph{Demostración.} Usando la definición de derivada,
\begin{align*}
    \frac{dC}{dt}
    &=\lim_{h\to 0}\frac{C(x^i(t+h),v^j(t+h),t+h)-C(x^i(t),v^j(t),t)}{h} \\
    &=\lim_{h\to 0}\frac{C(x^i(t),v^j(t),t)-C(x^i(t),v^j(t),t)}{h} \\
    &=\lim_{h\to 0}\frac{0}{h}
      =\lim_{h\to 0}0=0.
\end{align*}
Q. E. D.

Una propiedad importante de una constante del movimiento, además de ser constante, es que depende de la velocidad y la posición del sistema y no depende de la aceleración del sistema, por lo tanto, una constante del movimiento es una primera integral a las ecuaciones del movimiento.

Una constante del movimiento permite despejar una velocidad en función de las coordenadas y del resto de las otras velocidades,
\begin{equation*}
    C(x^i,v^j,t)=\text{constante}\quad\Longrightarrow\quad
    v^1=\frac{dx^1}{dt}=f(x^j,v^2,v^3,C).
\end{equation*}
Para el caso cuando la fuerza es conservativa, la velocidad se transforma en un factor integrando de las ecuaciones del movimiento que permite obtener la energía como una constante del movimiento.

Si se puede encontrar un número de constantes del movimiento igual al número de las coordenadas, o sea tres constantes del movimiento $\{C_1,C_2,C_3\}$, entonces tendríamos tres ecuaciones diferenciales de primer orden,
\begin{equation*}
\left\{
\begin{aligned}
    C_1&=C_1(x^j,\dot x^j),\\
    C_2&=C_2(x^j,\dot x^j),\\
    C_3&=C_3(x^j,\dot x^j),
\end{aligned}
\right.
\quad\Longrightarrow\quad
\left\{
\begin{aligned}
    \frac{dx^1}{dt}&=f^1(x^j,C_1,C_2,C_3),\\
    \frac{dx^2}{dt}&=f^2(x^j,C_1,C_2,C_3),\\
    \frac{dx^3}{dt}&=f^3(x^j,C_1,C_2,C_3).
\end{aligned}
\right.
\end{equation*}

\subsection{Ejemplo: sistema resorte--masa}

La acción de un resorte, con constante elástica $k$, sobre un cuerpo de masa $m$ genera la ecuación del movimiento
\begin{equation*}
    m\ddot x=-kx.
\end{equation*}
La fuerza elástica $F=(-)kx$, el signo $(-)$ indica que se opone a la fuerza aplicada, proviene del potencial
\begin{equation*}
    U(x)=\frac{1}{2}kx^2.
\end{equation*}
El punto de referencia de la energía elástica es $x_r=0$ donde $U(0)=0$.

La solución de la ecuación del movimiento es
\begin{equation*}
    x(t)=A\sin\left(\sqrt{\frac{k}{m}}t+\alpha\right),
\end{equation*}
donde hay dos constante del movimiento $A$ y $\alpha$. Estas constantes de integración pueden ser evaluadas conociendo las condiciones iniciales
\begin{equation*}
    x(0)=x_0,
    \qquad
    \dot x(0)=v_0.
\end{equation*}
Podemos relacionar las constantes de integración $A$ y $\alpha$ en función de las condiciones iniciales $x_0$ y $v_0$.

El sistema resorte-masa tiene la energía como constante del movimiento,
\begin{equation*}
    E=\frac{1}{2}m\dot x^2+\frac{1}{2}kx^2,
\end{equation*}
constante que depende de la posición y de la velocidad. Despejando la velocidad se tiene
\begin{equation*}
    \frac{dx}{dt}=\sqrt{\frac{2E}{m}-\frac{k}{m}x^2}.
\end{equation*}
Integrando se tiene
\begin{equation*}
    x(t)=\sqrt{\frac{2E}{k}}\sin\left(\sqrt{\frac{k}{m}}t+\alpha\right),
\end{equation*}
donde hay una sola constante de integración, $\alpha$, pues la constante $A$ queda expresada en función de la energía $E$.

\end{document}
